% 3. Design (1 page): design philosophy, considerations and constraints, functionalities of each simulator component focusing on the client-side simulator

\subsection{Philosophy}
The design philosophy for the client-side simulator is centered around modularity, simplicity, and extensibility \cite{6654207}. The goal is to create a robust and flexible simulation environment that can easily adapt to varying requirements and scenarios. Key principles include modularity, where each component of the simulator is designed as an independent module, allowing for easy updates, maintenance, and testing \cite{7463700}. This is achieved through the separation of functional components AUTH and PROCESS in order to maintain a modular code base. Simplicity is also crucial, as the architecture and codebase are kept as straightforward as possible to ensure readability and ease of understanding for future developers, this extends into a culture of detailed documentation throughout the project. \cite{9711975} Lastly, extensibility is emphasized to ensure the system can accommodate new features and functionalities with minimal changes to existing components which can be seen in FC logic whereby there is room to expand into further algorithms.

In terms of constraints and considerations, a few key details are addressed. This includes functionality and performance which is addressed by utilising Macquarie ASH servers to perform tests as appose to local infrastructure. This off-site approach removes the limitations of using personal devices as the infrastructure contains the necessary components required to perform the operations of DS-SIM and GreetClient.

Further to this, Reliability and usability is also key. Macquarie university ASH servers are maintained and scaled by Macquarie and are assumed to be reliable in terms of performance, latency and compatibility, ensuring accuracy and consistency of simulation results. With regards to usability, the user interface and interaction mechanisms must be implemented in a way to be user-friendly. This was achieved through Visual Studio Code remote host connection whereby the developer in this case was familiar with the layout and controls of the application despite being used across locations. This ensured that the application was both usable and accessible at any given time. Finally, It's important to highlight network considerations as a factor in the philosophy of development both at ASH server side and Locally, ensuring that regular code changes are committed to Github and ASH servers are available requires a stable network connection.


GreetClient at first was designed to handle the input and output data for the simulation, processing user inputs, sending commands to the server, and handling server responses. However after rounds of review, it was determined that GreetClient serves to be run via script without necessary Input and output from User, no user interface and no scale able requirements outside of meeting the script objectives. Therefore a design shift was decided to move away from component based parsers, error handlers and other user based interactive components and more into smooth functional components in preparation for automation. Design decisions around GreetClient kept the focus around being optimised for quick data processing and minimal latency by removing as many debug statements and unecessary components as possible


Another key consideration involves the logging capabilities of the GreetClient along with the error handling process. Effective logging is crucial for maintaining a reliable and debuggable system. Logs should be detailed enough to assist in debugging and troubleshooting, capturing essential information such as timestamps, event types, and relevant data points. This level of detail allows developers to trace the sequence of events leading up to an issue, making it easier to identify and resolve bugs. Although this particular implementation does not consider detailed logging capabilities, the initial design had this in place.

In addition to logging, detailed and intuitive error handling is necessary to ensure the simulator remains user-friendly and resilient. Error messages should be clear and informative, providing users with enough context to understand what went wrong and how to address it. This includes specifying the nature of the error, potential causes, and possible solutions or workarounds. User-friendly error handling also involves gracefully managing exceptions, ensuring that the simulator can recover from errors without crashing or losing data.

\subsection{DS-SIM Principles} 
The initial design for the project required an understanding of the ds-sim relationship with the client application. Scoping out and understanding each message between the client and server was critical to understanding and interpreting the application. The requirement was for a First Capable implementation of the ds-client written in our own words, and if possible to improve on this implementation

\subsection{Algorithm Architecture}
\subsubsection{First Capable}
At a high level, First Capable simply selects the first capable server that appears after requesting a server list. This results in a fairly resource intensive operation as in many cases, the server that performs the job may be over equipped in terms of memory, cpu and storage. First capable does not consider ready jobs, waiting jobs which makes it a potential crutch in large scale operations whereby the logic in server selection fails to take into account these considerations.The design principle of first Capable is outlined under implementation but at a high level, it essentially selects the first server capable of performing the job after the server list is given to the client.

\subsubsection{Out performing Algorithm considerations}
Although not implemented in this project, a range of alternative algorithms were tested in order to achieve innovation marks. The design aspects of these algorithms relied on the components presented, including memory, cpu, disk, state, waiting jobs and ready jobs. Considerations in this space involved worst fit scenario, whereby the server with the lowest spec was selected to perform the job, Best fit whereby the best server was selected for the job and a mixture of waiting job list logic. Finally, a rating system was tested in which each server is given a value based on a set of criteria. 
